\section{基于运行生产资料的汽轮机系统故障知识图谱构建}
\label{sec:knowledge_graph_v4}

\subsection{引言}

燃煤机组汽轮机系统由超高压缸、高压缸、中压缸、低压缸、再热系统、抽汽回热系统、给水系统及凝汽器冷端等多个设备共同组成，各设备之间通过蒸汽、凝结水和给水形成连续的物质与能量传递关系。第二章建立的状态监测模型能够在给定运行边界下获得设备健康状态参考，但状态偏差本身主要反映“当前运行状态相对健康基准发生了多大变化”，并不能直接说明异常参数属于哪个设备、多个异常之间是否具有实际热力联系，以及异常组合可能对应何种故障机理。因此，在利用状态偏差开展故障分类之前，还需要建立能够表征汽轮机设备拓扑、测点归属及故障传播规律的结构化知识模型。

知识图谱（Knowledge Graph, KG）以实体和关系为基本单元，将分散于系统结构图、DCS测点、运行规程、缺陷记录、检修资料和故障文献中的知识组织为可计算图结构\cite{Hogan2021KnowledgeGraphs}。汽轮机系统具有较明确的设备边界和介质传播路径，故障影响也往往沿这些结构逐级传播。Zhang和Hu针对汽轮机系统构建过程图与故障传播图，表明汽封缺陷、蒸汽泄漏及汽耗增加等故障信息能够依据设备机理形成层次化传播关系\cite{ZhangHu2021FaultPropagationGNN}。由此可见，将运行生产资料中的设备、参数和故障知识转化为知识图谱，可为后续图卷积网络提供区别于纯统计相关性的工程结构先验。

本章参照课题组已有的运行生产资料知识抽取与知识图谱构建技术路线\cite{YangHaodong2026CFB}，针对汽轮机系统重新定义知识范围、实体类型和关系类型。首先根据本机实际主通流、再热、抽汽回热和冷端结构分析参数关联与典型故障特性；随后利用科研智能体/大语言模型从多源资料中生成候选实体和三元组，并通过设备拓扑规则、实体类型规则、关系方向规则、语义归一和来源证据审核进行校验；最后形成经证据分级的汽轮机系统故障知识图谱，并建立领域参数节点与实际DCS测点之间的映射，为下一章KG-GCN模型输入构建提供基础。

\subsection{燃煤机组汽轮机系统参数及故障特性}
\label{subsec:turbine_fault_characteristics_v4}

\subsubsection{系统主通流及抽汽回热关系}

本文研究对象的主通流按照“主蒸汽--VHP--一次再热RH1--HP--二次再热RH2--IP--LP--凝汽器”的方向完成蒸汽膨胀做功与能量传递。主蒸汽压力、温度和流量决定VHP入口热力状态，VHP排汽形成一次冷再边界；一次高温再热参数进一步形成HP入口状态，HP排汽进入二次再热系统后形成IP入口边界，IP排汽继续进入低压缸，最终由凝汽器压力建立低压侧排汽边界。上述主通流关系均可由本机设备台账、DCS测点体系及第二章设备边界定义相互校核。

除主通流外，本机还具有多级抽汽回热支路。根据前期设备IO核对结果和最终测点体系，VHP设有一级抽汽并向1号高压加热器供汽；HP设有二级抽汽并与2号高压加热器相连；IP侧包含三级和四级抽汽，分别向3号和4号高压加热器供汽；LP侧包含六级、七级和八级抽汽，分别与6号、7号和8号低压加热器相连。设备台账还记录五抽、七抽至汽动给水泵小机的供汽支路。由此可见，抽汽并不是汽轮机本体之外的附加信息，而是汽缸状态向回热系统传播的重要介质通道。

需要指出的是，现有资料中关于除氧器蒸汽来源存在一处待最终核对的文字差异：最终测点体系将$10LAA10CP101$描述为五级抽汽/除氧器压力，而较早设备台账中出现“除氧器入口为四抽+凝结水”的记录。对于此类资料冲突，本文知识抽取过程不采用大语言模型直接裁决，而将相应关系标记为待核关系，在最终系统图或DCS证据确认前不进入KG-GCN参数图。该处理方式可避免将不同版本工程资料中的不一致信息误写为确定知识。

\begin{table}[htbp]
    \centering
    \caption{汽轮机系统主要设备及本机已确认抽汽回热关系}
    \label{tab:turbine_unit_topology_kg}
    \begin{tabular}{llll}
        \toprule
        汽轮机设备 & 抽汽级次 & 主要下游设备 & 证据来源 \\
        \midrule
        VHP & 一级抽汽 & 1号高压加热器 & 本机DCS/设备IO \\
        HP  & 二级抽汽 & 2号高压加热器 & 本机最终测点体系 \\
        IP  & 三级抽汽 & 3号高压加热器 & 本机最终测点体系 \\
        IP  & 四级抽汽 & 4号高压加热器 & 本机最终测点体系 \\
        LP  & 六级抽汽 & 6号低压加热器 & 本机最终测点体系 \\
        LP  & 七级抽汽 & 7号低压加热器 & 本机最终测点体系 \\
        LP  & 八级抽汽 & 8号低压加热器 & 本机最终测点体系 \\
        \bottomrule
    \end{tabular}
\end{table}

\subsubsection{典型通流及热力故障特性}

本文知识图谱重点描述能够通过DCS压力、温度、流量、阀位及冷端参数反映的通流和热力故障，不将需要专用高频振动或轴位移信号识别的纯机械故障作为主要对象。结合汽轮机故障文献、运行事件资料及本机测点条件，核心故障知识主要包括调节阀异常、通流部件侵蚀和沉积、汽封磨损与蒸汽泄漏、凝汽器换热性能劣化以及抽汽回热系统异常等。

（1）\textbf{调节阀死区及配汽异常。} 汽轮机调节阀通过改变有效流通面积调节进入汽轮机的蒸汽流量。Zhang和Hu针对汽轮机控制阀建立图模型，指出死区会使控制指令变化不能及时形成实际阀位响应，并可能诱发高压侧及系统振荡\cite{ZhangHu2021ControlValveGraph}。因此，可将“调节阀死区--阀位对指令不敏感--主蒸汽流量变化/系统振荡”作为直接故障知识链。对于某一具体VHP级组压力或温度的变化方向，则仍需根据机组工况和独立热力模型判断，不将其直接写入核心知识图谱。

（2）\textbf{汽封缺陷与蒸汽泄漏。} 级间汽封和轴端汽封的作用是限制不同压力区域之间的蒸汽泄漏。汽封齿磨损、间隙增大或结构损伤会使泄漏增加。汽轮机故障传播研究表明，HP、IP和LP级间或轴端汽封缺陷可进一步表现为蒸汽泄漏和汽耗增加，低压轴端汽封异常还可能引起空气进入低压系统并造成凝汽器真空下降\cite{ZhangHu2021FaultPropagationGNN}。针对600~MW超临界汽轮机HP--IP中间汽封泄漏的研究进一步指出，该泄漏会影响IP效率，热再温度和IP排汽温度是泄漏量估计的重要测量参数\cite{Xu2011HPIPLeakage}。因此，汽封故障知识采用“部件缺陷--蒸汽泄漏--性能异常/关联测量”的分层表达，而不直接假定所有温度和压力节点均具有固定方向响应。

（3）\textbf{通流部件侵蚀、沉积及表面退化。} Kubiak等对多台汽轮机运行诊断和检修结果进行总结，指出叶片表面粗糙度增加、固体颗粒侵蚀、迷宫或径向汽封磨损以及叶片沉积属于常见通流性能退化形式，关键压力与简化流量关系可用于运行阶段识别部件退化\cite{Kubiak2002TurbineDegradation}。其中固体颗粒侵蚀可能增大首级喷嘴有效面积；叶片或喷嘴沉积则会减小有效通流面积并降低输出能力和效率。关于沉积的数值研究进一步表明，通流面积减小会改变首级前后的压力分布，可利用压力变化估计堵塞程度\cite{Kubiak2002ScaleDeposit}。由于不同机型、级组和负荷条件下具体测点的变化方向可能不同，本文核心KG只保存“侵蚀--有效面积增大”“沉积--通流面积减小--效率/出力下降”等具有直接证据的机理关系；本机24个VHP状态点的点级响应由第五章半物理故障场景单独建模。

（4）\textbf{凝汽器污垢与冷端性能劣化。} 凝汽器水侧污垢会增加传热热阻并降低换热能力。针对燃煤电站凝汽器的在线研究表明，水侧污垢能够通过热有效性或污垢热阻进行监测，且冷端状态变化会进一步反映到凝汽压力和汽轮机低压侧运行状态\cite{Li2016CondenserFouling,Medica2018CondenserCondition}。因此，图谱中保存“凝汽器水侧污垢--换热能力下降--凝汽器背压受影响”的知识链，并将循环水入口温度和流量作为冷端外部边界参数，而不是将其下降或升高本身定义为设备故障。

（5）\textbf{抽汽回热及给水加热器异常。} 给水加热器内部管束、管板或隔板泄漏会破坏正常抽汽--给水换热过程并降低回热性能。Heo和Lee针对闭式给水加热器内部泄漏开展诊断研究，指出内部泄漏会影响整个汽轮机循环性能\cite{Heo2012FWHLeakage}。ASME TDP-1还将汽轮机抽汽系统、给水加热器、疏水以及防止水或湿蒸汽倒灌列为汽轮机水损伤防护的重要对象。因此，抽汽逆止阀、加热器泄漏及抽汽管路异常可以作为系统级故障知识保存在领域图谱中；但只有与本文机组结构和DCS可观测量兼容的关系才参与后续实时诊断。

\subsection{基于运行生产资料的汽轮机故障知识图谱构建}
\label{subsec:turbine_kg_construction_v4}

\subsubsection{运行生产资料及知识来源}

知识图谱的可靠性首先取决于知识来源。本文所使用的运行生产资料主要分为四类。第一类为本文机组的系统结构资料和设备台账，用于确认VHP、HP、IP、LP、再热系统、凝汽器及各级抽汽回热设备之间的实际物理连接；第二类为DCS测点与设备IO资料，用于确定领域参数概念与实际测量点之间的对应关系；第三类为运行规程、缺陷与检修记录等工程资料，用于补充故障发生部位、现象和处置经验；第四类为同行评议论文、行业标准和公开运行事件报告，用于补充具有可追溯来源的通流故障机理和跨设备传播知识。

与通用汽轮机知识不同，本机设备拓扑优先采用本文项目已逐点整理的系统台账和DCS资料。例如，一级抽汽压力、一级抽汽温度及1号高加入口抽汽参数均可由实际点号直接确认；二至八级抽汽中的主要压力和温度概念也可由最终测点体系映射。对于外部论文，则主要抽取“故障--异常机理--性能影响”关系，而不直接替代本文机组的设备连接关系。通过这种方式，将“本机结构事实”和“外部故障机理知识”分层融合，避免将其他机型的抽汽结构或测点布置照搬到本文对象。

\subsubsection{实体与关系模式设计}

结合后续故障诊断需求，本文将核心实体划分为设备、参数、异常和故障四类，并保留处置实体作为工程解释扩展。设备实体用于表示VHP、HP、IP、LP、再热器、各级抽汽系统、加热器和凝汽器等物理对象；参数实体表示能够对应实际DCS数值的压力、温度、流量和阀位等变量；异常实体表示“蒸汽泄漏”“换热能力下降”“通流面积减小”等状态变化；故障实体表示“汽封缺陷”“固体颗粒侵蚀”“叶片沉积”等根因或部件失效模式。

知识关系根据工程语义划分为结构关系和故障关系两组。结构关系包括“介质流向”“抽汽至”“供汽至”“归属”和“构成边界”等，用于表达汽轮机系统真实物理拓扑；故障关系包括“发生于”“导致”“影响”“表现为”和“关联监测”等，用于描述故障从部件到异常状态再到监测参数的传播链。任一三元组记为
\begin{equation}
    \tau=(h,r,t),
\end{equation}
其中$h$和$t$分别为头实体和尾实体，$r$为关系类型。通过统一实体和关系模式，可将原本分散在表格、文字和系统图中的知识转换为统一图结构。

\begin{table}[htbp]
    \centering
    \caption{汽轮机故障知识图谱主要实体及关系类型}
    \label{tab:turbine_kg_schema_v4}
    \begin{tabular}{lll}
        \toprule
        类别 & 类型 & 示例 \\
        \midrule
        实体 & 设备 & VHP、一级抽汽、1号高加、凝汽器 \\
        实体 & 参数 & VHP排汽温度、一级抽汽压力、凝汽器背压 \\
        实体 & 异常 & 蒸汽泄漏、通流面积减小、换热能力下降 \\
        实体 & 故障 & 固体颗粒侵蚀、叶片沉积、汽封缺陷 \\
        关系 & 结构关系 & 介质流向、抽汽至、供汽至、归属 \\
        关系 & 故障关系 & 发生于、导致、影响、表现为、关联监测 \\
        \bottomrule
    \end{tabular}
\end{table}

\subsubsection{科研智能体辅助知识抽取与实体标准化}

针对运行规程、故障分析和检修文本中非结构化知识量较大的问题，本文采用科研智能体/大语言模型辅助生成候选三元组。首先按照预定义实体和关系Schema对原始段落进行解析，识别设备名称、运行参数、异常现象、故障名称及处置动作，并生成候选$(h,r,t)$；随后利用规则程序检查实体类型和关系方向，例如“参数--归属--设备”“故障--导致--异常”等组合必须满足预定义类型约束，不满足约束的候选关系直接删除或退回人工复核。

不同资料对同一实体往往存在不同表述，例如“HP--IP中间汽封泄漏”“N2 packing leakage”和“中间汽封漏汽”可能指向同一机理对象。对此，先通过设备位置、上下游对象和关键词进行规则归一，再使用文本嵌入模型计算候选名称的语义相似度\cite{Reimers2019SentenceBERT}。对于高相似候选，由系统给出合并建议；最终标准名称仍以本文设备命名体系和工程资料为准。该过程使大语言模型承担信息检索和候选生成任务，而不直接拥有“创造知识”的权限。

\subsubsection{证据分级与知识校验}

为保证图谱关系可追溯，本文在候选知识进入核心语义层前设置证据分级。A0级表示由本文机组设备台账、DCS测点或冻结IO直接确认的结构事实；A1级表示由同行评议原始论文、行业标准或官方运行资料直接支持的故障机理关系。对于不同资料之间仍存在冲突的本机关系，标记为A0$^*$并暂不用于模型拓扑；仅依据热力机理或大语言模型推断得到、尚缺乏直接点级证据的关系标记为B级。

知识校验主要包括四方面：一是\textbf{实体类型校验}，实时参数节点必须能够表示一个可测物理量，诸如“背压升高”“阀位异常”等描述只能作为异常节点；二是\textbf{设备拓扑校验}，本机结构边必须与设备IO和介质流向一致；三是\textbf{故障证据校验}，当文献只证明“汽封缺陷导致蒸汽泄漏”时，不进一步无依据地扩展为“汽封缺陷必然导致某一具体DCS温度升高”；四是\textbf{DCS映射校验}，用于后续GCN的参数节点必须能够映射本文机组实际测点。

在这一规则下，知识抽取形成“候选关系生成--Schema规则校验--实体标准化--证据审核--图谱入库”的闭环。科研智能体和知识图谱还可以反向进行一致性检查：若一个参数被同时归属于两个不相容设备、一个故障关系方向与已有物理链冲突，或某个所谓参数节点无法对应实际测量值，则自动标记为待复核项，从而降低大规模文本抽取时的错误传播。

\subsubsection{知识图谱构建结果}

按照上述流程对本机设备拓扑、DCS测点及公开故障知识进行抽取和审核，形成汽轮机系统故障知识图谱核心层。经第二轮拓扑与证据审计后，当前核心图谱包含127个节点和123条A0/A1主关系，其中81条为本机A0结构/参数事实边，40条为外部A1直接故障机理边，另有2条本机资料冲突关系以A0$^*$形式保留但不参与当前模型投影。与此同时，共形成45条领域参数--DCS测点映射，其中44条已确认；7条缺乏直接点级证据但具有热力合理性的关系单独保存在B级推断层，不进入主知识图谱。

从节点结构看，知识图谱同时覆盖主通流、再热、抽汽回热、给水和凝汽器冷端。与第一版通用化抽汽拓扑相比，审计后的结构恢复了VHP、HP、IP和LP对应的多级抽汽支路，使系统知识能够反映本文机组真实设备接口。故障知识部分则形成调节阀异常、汽封泄漏、通流部件侵蚀与沉积、凝汽器性能劣化和回热设备异常等多个故障子图。对于主图中每一条边，均保存其来源类型、兼容范围和是否允许参与GCN投影等属性，便于后续模型使用和故障解释。

\begin{figure}[htbp]
    \centering
    % 最终矢量图由 turbine_kg_v2_audited_main_edges.csv 自动生成，待SVG纳入仓库后替换
    \fbox{\parbox[c][6.0cm][c]{0.94\textwidth}{\centering
    汽轮机系统故障知识图谱核心层：\\
    上部为本机主通流与1/2/3/4/6/7/8级抽汽回热拓扑；\\
    中部为实际DCS参数概念映射；\\
    下部为经A1证据审核的调节阀、汽封、通流、凝汽器及回热故障机理子图。}}
    \caption{汽轮机系统故障诊断知识图谱核心层}
    \label{fig:turbine_kg_v4}
\end{figure}

需要说明的是，图~\ref{fig:turbine_kg_v4}用于展示核心结构，并未将127个节点全部展开。完整节点和关系表在模型实现阶段以结构化文件保存。知识图谱的系统级覆盖范围并不意味着本文已经对所有设备完成故障分类实验；第五章仅选择VHP作为代表性对象，通过领域参数--DCS映射从全系统知识图谱中提取相应参数子图，用于验证本文KG-GCN诊断方法。

\subsection{本章小结}

本章面向燃煤机组汽轮机系统故障诊断中的结构知识表达问题，构建了基于运行生产资料的汽轮机故障知识图谱。首先依据本机主通流和多级抽汽回热结构分析参数关联及故障传播特性，并结合公开汽轮机故障研究整理调节阀异常、汽封泄漏、通流部件侵蚀与沉积、凝汽器污垢及回热设备异常等典型故障知识；随后定义设备、参数、异常、故障等实体及相应关系模式，采用科研智能体完成候选知识抽取，并通过Schema规则、语义标准化和来源证据进行校验。经审计形成的核心语义层包含127个节点和123条主关系，同时建立领域参数与实际DCS测点之间的映射。所构建知识图谱既保存汽轮机系统真实物理拓扑，又保存具有直接证据的故障传播知识，为下一章将状态监测残差加载至参数节点并开展图卷积故障诊断提供结构基础。
